\documentclass[10pt]{article}

\usepackage[utf8]{inputenc}

\usepackage[T1]{fontenc}

\usepackage{amsmath}

\usepackage{amsfonts}

\usepackage{amssymb}

\usepackage[version=4]{mhchem}

\usepackage{stmaryrd}

\usepackage{mathrsfs}

\usepackage{bbold}



\begin{document}

Метод ренормализационной группы был развит в середине 50-х гг. Э. Штюкельбергом (E. Stueckelberg) и А. Петерманом (A. Peterman), М. Гелл-Маном (M. Gell-Mann) и Ф. Лоy (F. Low), Н. Н. Боголюбовым и Д. В. Ширковым. Помимо квантовой теории поля этот метод стал успешно применяться в статистич. физике, а также в нек-рых разделах физики, изучающих не только микроскопич., но и макроскопич. процессы.



Лит.: [1] Боголюбов Н. Н., Ши рков Д. В., Введение в теорию квантованных полей, 4 изд., М., 1984; [2] За в ьяло в О. И., Перенормированные диаграммы Фейнмана, М., 1979; [3] И ц и кс он К., 3 ю бер Ж.-Б., Квантовая теория поля, пер. с англ., М., 1984; [4] Колл и нз Дж., Перенормировка, пер. с англ., М., 1988; [5] С м и р н о в В. А., Перенормировка и асимптотические разложения фейнмановских амплитуд, М., $1990 . \quad B . A$. Смирнов. ПЕРЕНОРМИРУЕМОСТЬ - см. Глэшоу-Вайнберга-Салама модель, Квантовая теория гравитации.



ПЕРЕНОСА УРАВНЕНИЕ БОЛЬЦМАНА - СМ. БоЛЬЩМана кинетическое уравнение.



ПЕРЕОПРЕДЕЛЕННАЯ СИСТЕМА - см. Дифференциальное уравнение с частными производными второго порядка.



ПЕРЕПЛЕТЕНИЯ ОПЕРАТОP - см. Сплетающий оператор. ПЕРЕПОЛНЕННАЯ СИСТЕМА - семеЙствО векТОрОВ $\left\{e_{x} ; x \in \mathscr{Y}\right\}$ в гильбертовом пространстве $H$, удовлетворяющее свойству полноты



\[

\|\psi\|^{2}=\int_{\mathscr{X}}\left(e_{x}, \psi\right)^{2} \mu(d x) \text { для всех } \psi \in H,

\]



где $\mu$ - нек-рая мера на $\mathscr{X}$. Всякий вектор $\psi \in H$ допускает разложение по векторам П. с.



\[

\psi=\int_{x} e_{x}\left(e_{x}, \psi\right) \mu(d x)

\]



соответствующее представление имеет место и для операторов в $H$. Частным случаем П. с. является полная ортогональная система в $H$, однако в общем случае векторы $e_{x}$ могут быть неортогональными и даже линейно зависимыми. Важнейший пример П. с.- система когерентных состояний.



Jum.: [1] Б е р е з и н Ф.А., Метод вторичного квантования, 2 изд., М., 1986; [2] Клауде р Дж., Сударш ан Э., Основы кванговой оптики, пер. с англ., М., $1970 . \quad$ А. С. Холево.



ПЕРЕСЕЧЕНИЙ ФОРМА - билинейная форма (.०.) на пространстве гомологий $H_{n}(M) 2 n$-мерного ориентированного многообразия $M$, определяемая по следующему правипу. Пусть $\alpha$ и $\beta$ - ориентированные цепи, представляющие два цикла $[\alpha]$ и $[\beta]$ в $H_{n}(M)$. Каждой точке пересечения $\alpha$ и $\beta$ ставится в соответствие знак \textbackslash pm 1 , в зависимости от того, индуцируют ли ориентации $\alpha$ и $\beta$ ориентацию $M$ или противоположную. Тогда $(\alpha \circ \beta)$ равно сумме всех знаяов.



С. В. Чмутов.



ПEPECTAHОВОЧНЫЕ СООТНОШЕНИЯ - соотнОШения, описывающие различие в воздействии на состояния квантовой системы последовательно применяемых операгоров, взятых в разном порядке, то есть отличие $A B|\psi\rangle$ от $B A|\psi\rangle$. В математич. аппарате квантовой теории П. с. отражают свойство некоммутативности наблюдаемых алгебры, әвязанное с сушествованием физич. величин, не допускающих в принципе одновременное измерение. Отличие от нуля коммутатора $[A, B] \div A B-B A$ операторов, описывающих эти величины, приводит к сушествованию соотношений неопределенности. Коммутатор является квантовым дналогом классич. скобок Пуассона и приводит к наличию гтруктуры Ли в алгебре наблюдаемых. Примерами некомиутирующих величин являются операторы координаты и импульса, компоненты оператора углового момента, компоненты спина и т. п. Особый класс образуют П. с. генерагоров динамич. групп Ли и генераторов групп, описываюших свойства симметрии рассматриваемой физич. системы (напр., генераторы групп Галилея и Пуанкаре или групп унитарной симметрии $S U(n)$ и т.п.). Важное значение в квантовой теории играют П.с. токов алгебр. В последние годы в связи с появлением концепции суперсимметрии, квантованием нелинейных уравнений и развитием теории квантовых струн активно используются П.с. алгебр Ли бесконечномерных групп Ли, таких, как группа Вирасоро или группа Каца - Муди. Фундаментальное значение имеют П. с. канонич. переменных (или полей), соответствующих бозевским или фермиевским степеням свободы, и П. с. фундаментальных токов (модель нерелятивистских токов, модель Сугавара и т. Д.).



П.с. для бозевских степеней свободы. В квантовой механике систем с конечным числом $d<\infty$ степеней свободы постулируются П. с. вида



\[

\begin{aligned}

& {\left[Q_{j}, P_{k}\right]=i \delta_{j k} \mathbb{1},\left[Q_{j}, Q_{k}\right]=0,\left[P_{j}, P_{k}\right]=0,} \\

& j, k=1, \ldots, d, \hbar=1 .

\end{aligned}

\]



Здесь $Q_{j}$ - обобщенные координаты, а $P_{j}$ - канонически сопряженные импульсы. П.с. (1) обычно называют кано ническими коммутационными соотношения ми в форме Гей зен берга. Известно, что соотношения (1) нельзя удовлетворить ограниченными операторами и они выполняются только на плотной области в гильбертовом пространстве $H$. От канонич. операторов $Q$ и $P$ часто переходят к унигарным операторам Вейля $U(p)=e^{i p Q}, V(q)=e^{i q P}$, $q, p \in \mathbb{R}$, к-рые удовлетворяют в силу (1) ком мутац ионным соотношениям в форме Вейля:



\[

\left.\begin{array}{l}

V(q) U(p)=e^{i p q} U(p) V(q), \\

U\left(p_{1}\right) U\left(p_{2}\right)=U\left(p_{1}+p_{2}\right), V\left(q_{1}\right) V\left(q_{2}\right)=V\left(q_{1}+q_{2}\right)

\end{array}\right\}

\]



(для упрощения обозначений рассматривается случай $d=1$ ). Соотношения (2) определяют группу Гейзенберга - Вейля, алгеброй Ли к-рой служит алгебра Гейзенберга, порождаемая П.с. (1). Наряду с канонич. операторами $\boldsymbol{Q}$ и $P$ часто рассматривают операторы рождения и гибели процессов



\[

a^{+}=\frac{1}{\sqrt{2}}(Q+i P), \quad a=\frac{1}{\sqrt{2}}(Q-i P) .

\]



Это замкнутые неограниченные операторы, удовлетворяюшие в силу (1) П. с. вида



\[

\left[a, a^{+}\right]=1,[a, a]=0,\left[a^{+}, a^{+}\right]=0,

\]



к-рые также выполняются только на плотной в $H$ области.



В релятивистской квантовой теории поля постулируются П. с. полевых операторов в разных точках. Напр., для нейтрального скалярного поля $\varphi(x)$ массы $m$ П. с. имеют вид:



\[

[\varphi(x), \varphi(y)]=-i \Delta(x-y),

\]



где $\Delta-$ инвариантная причинная функция, к-рая является решением уравнения $\left(\square-m^{2}\right) \Delta=0$ и обращается в нуль при пространственноподобном аргументе. В канонич. подходе используют одновременные П. с. между полем $\varphi(x, t)$ и канонически сопряженным импульсом $\pi(x, t)=\frac{\delta z}{\delta \dot{\varphi}}$ (точ-

нее плотность импульса) вида:



\[

\left.\begin{array}{l}

{[\varphi(x, t), \pi(y, t)]=i \delta(x-y),} \\

{[\varphi(x, t), \varphi(y, t)]=0,} \\

{[\pi(x, t), \pi(y, t)]=0 .}

\end{array}\right\}

\]



Поскольку $\varphi(\boldsymbol{x}, t), \pi(\boldsymbol{y}, t)$ являются операторно-значными распределениями, то для получения операторов производится сглаживание с функциями из подходящего действительного векторного пространства $V$ пробных функций, снабженного скалярным произведением



\[

\therefore \quad(f, g)=\int f(\boldsymbol{x}) g(\boldsymbol{x}) d^{3} \boldsymbol{x},

\]



то есть вводятся операторы



\[

\begin{aligned}

& \varphi(f)=\int \varphi(x, t) f(x) d^{3} x \\

& \pi(g)=\int \pi(x, t) g(x) d^{3} x

\end{aligned}

\]



удовлетворяющие в силу (4) гейзенберговским канонич. коммутационным соотношениям



$[\varphi(f), \pi(g)]=i(f, g),[\varphi(f), \varphi(g)]=0,[\pi(f), \pi(g)]=0$.



ПЕРЕСТАНОВОЧНЫЕ 433





\end{document}